\begin{solapartat}
\punts{0,1} Càlcul de la mitjana aritmètica:

\begin{center}
\begin{tabular}{lllll}
\toprule
\textbf{x (cm)} & $\mathbf{V_\mathrm{I}}$ \textbf{(V)} & $\mathbf{V_\mathrm{II}}$ \textbf{(V)} & $\mathbf{V_\mathrm{III}}$ \textbf{(V)} & $\mathbf{V_\mathrm{mitj\grave{a}}}$ \\
\midrule
4,00 & 1,4 & 1,5 & 1,4 & 1,43 \\
8,00 & 2,8 & 2,9 & 2,8 & 2,83 \\
12,0 & 4,2 & 4,3 & 4,4 & 4,30 \\
16,0 & 5,7 & 5,7 & 5,8 & 5,73 \\
\bottomrule
\end{tabular}
\end{center}

\punts{0,2} Les línies equipotencials són paral·leles a les plaques.

\punts{0,25} Les línies de camp són perpendiculars a les línies equipotencials i van dirigides del potencial alt al potencial baix.

\figura[0.55]{sol_fig1.png}

\destaca{Nota:} no cal que a la solució s'indiquin els valors del potencial.

\punts{0,4} Representació gràfica:

\figura[0.45]{sol_fig2.png}

No incloure el títol a l'eix resta \destaca{0,1 p}

No incloure les unitats a l'eix resta \destaca{0,2 p}

Cal observar que disposem d'un punt més que el donat a l'enunciat, si connectem el terminal negatiu a la placa connectada al terminal negatiu de la font d'alimentació i, a més, aquesta placa està ubicada a $x = 0$~cm, llavors en aquest punt el potencial elèctric és zero. No es penalitzarà l'omissió d'aquest punt.

No escalar o representar correctament la gràfica serà qualificat com a zero.

\punts{0,3} Ja hem vist que el camp elèctric és perpendicular a les plaques i també hem determinat el seu sentit, per tant el camp només té component en l'eix de les abscisses i el seu mòdul és igual al valor absolut de la component $x$ del camp elèctric. El camp és la derivada del potencial respecte a la posició o directament el pendent de la representació $V(x)$. Com la gràfica resultat s'ajusta molt bé a una recta, llavors el camp és constant, no depèn de la posició. A partir del pendent de la recta obtenim:
\[ E = 0,358\ \mathrm{V/cm} = 35,8\ \mathrm{V/m} \]
Idealment el pendent s'obté d'un ajust per mínims quadrats o a partir de representar la recta que millor s'ajusta als punts.

\destaca{Alternativament, també es considerarà correcte} que per determinar el pendent faci servir un parell de punts de la taula \destaca{i també es considerarà correcte} que l'estudiant apliqui directament $E = \frac{\Delta V}{\Delta x}$ sempre que justifiqui que aquesta equació és vàlida en el cas de camps constants, és a dir, quan tenim el camp creat per dues plaques planes i paral·leles.
\end{solapartat}

\begin{solapartat}
\punts{0,3} Seguirà una trajectòria paral·lela a les línies de camp i com que la càrrega és positiva tindrà el mateix sentit, per tant, es mourà en línia recta segons l'eix $x$ cap a la placa negativa.

\punts{0,35} La direcció serà, doncs, horitzontal i sentit cap a la placa negativa:

\figura[0.5]{sol_fig3.png}

\punts{0,3} El mòdul de la força és
\[ F = qE = 3\times 10^{-3}\times 35,8 = 0,107\ \mathrm{N} \]

\punts{0,3} I el treball és:
\[ W = -q\Delta V = -3\times 10^{-3}(0 - 2,83) = 8,49\times 10^{-3}\ \mathrm{J} \]
\end{solapartat}
